\documentclass[a4paper,conference]{IEEEtran}

% -----------------------------
% Packages
% -----------------------------
\usepackage{graphicx}
\usepackage{ragged2e}
\usepackage{amssymb}
\usepackage{amsmath}
\usepackage{booktabs}
\usepackage{array}
\usepackage{cite}
\usepackage{url}
\newcommand{\safeincludegraphics}[2][]{%
\IfFileExists{#2}{\includegraphics[#1]{#2}}{\fbox{\parbox[c][0.35\textheight][c]{0.92\linewidth}{\centering Missing figure: #2}}}}

% -----------------------------
% Custom Column Types
% -----------------------------
\newcolumntype{L}[1]{>{\raggedright\arraybackslash}p{#1}}
\newcolumntype{C}[1]{>{\centering\arraybackslash}p{#1}}
\newcolumntype{R}[1]{>{\raggedleft\arraybackslash}p{#1}}

% -----------------------------
% Title / Author
% -----------------------------
\title{Open-Set Qur'anic Reciter Recognition Using Hybrid Acoustic Features: A Comparative Evaluation of Random Forest and Bidirectional LSTM Models}

\author{
\IEEEauthorblockN{Abdusamii Abdulazeez\IEEEauthorrefmark{1} and Habeeb Omotunde\IEEEauthorrefmark{2}}
\IEEEauthorblockA{\IEEEauthorrefmark{1}Information Technology, Imam Mohammad Ibn Saud Islamic University, Riyadh, Saudi Arabia\\
443022312@sm.imamu.edu.sa}
\IEEEauthorblockA{\IEEEauthorrefmark{2}Information Systems, Imam Mohammad Ibn Saud Islamic University, Riyadh, Saudi Arabia\\
homotunde@imamu.edu.sa}
}

% -----------------------------
% Document Start
% -----------------------------
\begin{document}

\maketitle
\begin{abstract}
Automatic Qur'anic reciter recognition is relevant to digital learning and recitation verification, yet limited attention has been given to inputs from previously unseen reciters. This paper presents Qurr\={a}, a small-scale open-set proof of concept combining hybrid acoustic descriptors with a reliability-based rejection rule. A Random Forest (RF) using a 511-dimensional clip-level hybrid vector is compared with a Bidirectional LSTM (BLSTM) baseline using 13 clip-level MFCC summaries. The evaluation uses a 2,535-recording development corpus from five reciters and a content-separated external set of 399 recordings from seven reciters. Before rejection, RF and BLSTM achieved 100.0\% and 97.81\% accuracy, respectively, on 228 known-reciter recordings from unseen surahs. On 171 recordings from three unseen reciters, RF rejected 97.66\% and BLSTM rejected 87.13\%; their known-reciter false-rejection rates were 6.58\% and 9.21\%. Overall external accuracies after rejection were 95.24\% and 88.97\%. These results support hybrid-feature RF for this dataset, but do not isolate classifier architecture from input representation.
\end{abstract}

\begin{IEEEkeywords}
Qur'anic reciter recognition, open-set recognition, Random Forest, bidirectional LSTM, MFCC, acoustic features, speech processing, Arabic speech recognition, machine learning, speaker identification.
\end{IEEEkeywords}

% -----------------------------
% Main Text
% -----------------------------
\section{Introduction}

The Holy Qur'an has been preserved for centuries through authenticated oral transmission and precise recitation practices. Accurate recitation plays a central role in Qur'anic learning and religious practice, motivating computational systems capable of analyzing recitation audio for learning, verification, and accessibility. With the growth of digital learning platforms and artificial intelligence applications in education, there is increasing interest in automated Qur'anic audio analysis systems \cite{ref2, ref9, ref10}.

Automatic Qur'anic recitation analysis has attracted significant research attention in recent years. Existing studies have explored reciter recognition, speaker identification, recitation classification, and Arabic speech processing using classical machine learning, deep learning, and speech-recognition methods \cite{ref1, ref3, ref4, ref5, ref6, ref7, ref8, ref11, ref12}. However, most reciter-recognition studies use closed-set evaluation, where all test reciters are present during training. This does not reflect deployment conditions where systems may encounter previously unseen reciters.

This paper presents Qurr\={a}, an open-set Qur'anic reciter recognition framework combining MFCC, chroma, spectral, rhythmic, and tonal descriptors with reliability-based rejection. An RF classifier \cite{breiman} using the full hybrid representation is compared with a BLSTM baseline \cite{schuster} using only MFCC summaries. Model development uses Surahs 78--103, while external evaluation uses Surahs 104--114 and includes known and unseen reciters. This separation removes verse overlap, but does not control recording source or channel.

The main contributions of this work are as follows:

\begin{enumerate}
\item An open-set Qur'anic reciter recognition framework using hybrid acoustic feature representations and reliability-based rejection.
\item A comparison of hybrid-feature RF and an MFCC-summary BLSTM baseline under fixed data partitions.
\item A content-separated protocol that distinguishes raw known-reciter classification, known-reciter false rejection, and unseen-reciter rejection.
\item Per-reciter and aggregate evidence for a small-scale, domain-specific proof of concept.
\end{enumerate}

\section{Related Work}

Existing studies show that Qur'anic reciter identity can be inferred from audio using deep-learning models \cite{ref1,ref6,ref7} and conventional acoustic classifiers \cite{ref3,ref4}. Related systems address recitation-rule checking, memorization support, and verse-level processing \cite{ref2,ref9,ref10}. These works establish the feasibility of automated Qur'anic audio analysis, but they commonly assume that every test reciter belongs to a known training class. Open-set recognition addresses this limitation by allowing a model to reject samples outside the training classes \cite{openset}. This distinction is important for Qur'anic applications, where a user may submit audio from any reciter rather than only reciters represented in the training corpus.

Two gaps motivate the evaluation design. First, closed-set accuracy can overstate deployment performance because every input is forced into a known class. Second, Qur'anic recordings share fixed text, so verse overlap can reward content familiarity. We therefore separate surahs and introduce unseen reciters. This design tests content-separated behavior, although source and channel effects remain possible and are treated as threats to validity.

\section{Proposed System}

\subsection{System Overview}

The proposed Qurr\={a} framework is an open-set Qur'anic reciter recognition system designed to identify known reciters from audio recordings while rejecting previously unseen reciters. The system consists of four main stages: audio preprocessing, acoustic feature extraction, classification, and reliability-based decision making.

Given an input Qur'anic recitation clip, the system first applies standardized preprocessing to ensure uniform sampling and signal quality. Acoustic features are then extracted using a hybrid feature extraction pipeline. These features are passed to either RF or BLSTM classification, and a reliability-based module determines whether the prediction corresponds to a known reciter or should be rejected as unknown.

The overall architecture is illustrated in Fig.~\ref{figure1}.

% -----------------------------
% Figure 1
% -----------------------------
\begin{figure*}[t]
\centering
\safeincludegraphics[width=0.82\textwidth]{figure1.png}
\caption{Overall architecture of the proposed Qurr\={a} open-set reciter recognition framework.}
\label{figure1}
\end{figure*}

\subsection{Audio Preprocessing and Feature Extraction}

All recordings are loaded in full, converted to mono, noise-reduced, amplitude-normalized, silence-trimmed at 30~dB, and resampled to 22,050~Hz using librosa \cite{librosa}. No common-duration padding or truncation is applied. Frame-level MFCC, chroma, mel-spectrogram, spectral contrast, rolloff, centroid, tempogram, tonnetz, RMS, and zero-crossing descriptors are temporally mean-pooled and concatenated into a fixed 511-dimensional clip vector for RF. Thus, variable clip length is handled by aggregation.

The hybrid representation is designed to capture complementary properties of recitation audio. MFCC and mel-spectrogram features summarize cepstral and frequency-band structure, chroma and tonnetz features reflect tonal organization, spectral descriptors capture timbral differences, and rhythm-related features represent temporal delivery patterns. Since reciters may differ in pitch contour, articulation, pacing, and spectral color, combining these descriptors provides a broader acoustic profile than relying on MFCC features alone.

Table~\ref{table_feature_blocks} gives the exact dimensional composition implemented by the extractor. The first- and second-order delta operations are each reduced to one scalar after the 32 MFCC summaries have already been pooled. The 384-bin mean tempogram is consequently the largest block. Reporting this composition is important because ``hybrid features'' otherwise conceals a strongly unequal allocation of dimensions, and it makes clear that RF receives substantially more rhythm-related information than the 13-value BLSTM input.

\begin{table}[htbp]
\centering
\caption{Implemented composition of the 511-dimensional RF vector.}
\label{table_feature_blocks}
\renewcommand{\arraystretch}{1.04}
\setlength{\tabcolsep}{4pt}
\begin{tabular}{l r}
\toprule
\textbf{Feature block} & \textbf{Dimensions} \\
\midrule
MFCC summaries & 32 \\
Mean first-/second-order MFCC delta & 2 \\
Chroma and mel summaries & 12 + 64 \\
Spectral contrast, rolloff, centroid & 7 + 1 + 1 \\
Zero-crossing rate and RMS & 1 + 1 \\
Tempogram & 384 \\
Tonnetz & 6 \\
\midrule
Total & 511 \\
\bottomrule
\end{tabular}
\end{table}

% -----------------------------
% Figure 2
% -----------------------------
\begin{figure}[htbp]
\centering
\safeincludegraphics[width=\columnwidth]{figure2.png}
\caption{Hybrid acoustic feature extraction and processing pipeline.}
\label{figure2}
\end{figure}

\subsection{Random Forest Classifier}

The RF classifier uses 100 decision trees with maximum depth 10 and random state 42. It is trained on the full 511-dimensional hybrid feature vector, and probability calibration is applied to improve confidence estimates used by the open-set rejection stage. Predictions are aggregated from individual trees using majority voting, allowing the model to operate robustly on high-dimensional acoustic representations without requiring explicit temporal modeling.

\subsection{Bidirectional Long Short-Term Memory Classifier}

The BLSTM baseline uses the first 13 clip-level MFCC summary values from the aggregated vector. Each 13-dimensional vector is repeated across 16 input positions to meet the fixed BLSTM input shape; it is therefore not a frame-level temporal representation. The network contains one bidirectional LSTM layer with 64 hidden units per direction, attention \cite{bahdanau}, and a fully connected classifier. Training uses AdamW (learning rate 0.0005), batch size 64, dropout 0.5, early stopping, and learning-rate scheduling. This baseline evaluates the implemented pipeline, not a feature-matched architecture comparison.

\subsection{Reliability-Based Open-Set Recognition}

To enable open-set recognition, each prediction is accepted only when confidence and distance criteria indicate that the sample belongs to a known class. Otherwise, the output label is \textbf{Unknown}.

Let $p_1$ and $p_2$ denote the highest and second-highest predicted class probabilities, respectively. Let $d(x,c_1)$ denote the Euclidean distance between feature vector $x$ and the centroid of the predicted class $c_1$, and let $T(c_1)$ be the 95th-percentile intra-class distance threshold estimated from training data. The prediction is accepted only when
\begin{equation}
p_1 \geq 0.95 \ \land\ p_2 < 0.10 \ \land\ (p_1-p_2) \geq 0.80 \ \land\ \frac{d(x,c_1)}{T(c_1)} \leq 1.0.
\label{eq:reliability}
\end{equation}

This mechanism allows the system to operate in open-set conditions without explicitly training on an unknown class. The workflow is shown in Fig.~\ref{figure3}.

The 95th-percentile distance threshold is estimated from model-fitting samples. The probability, secondary-probability, margin, and distance-ratio cutoffs are manually selected engineering operating points rather than values optimized on the external set. The rule rejects ambiguous predictions and high-confidence samples lying outside the predicted class region. No threshold sensitivity study was performed.

% -----------------------------
% Figure 3
% -----------------------------
\begin{figure*}[t]
\centering
\safeincludegraphics[width=0.72\textwidth]{figure3.png}
\caption{Reliability-based open-set recognition decision workflow.}
\label{figure3}
\end{figure*}

\section{Experimental Setup}

\subsection{Dataset}

Experiments use recordings downloaded from reciter-specific EveryAyah archives. The 2,535-recording development corpus is split into 2,028 model-fitting and 507 holdout recordings from Abdulbasit Abdulsamad, Abdulrahman Alsudaes, Ahmad Al-Ajmy, Maher Al Meaqli, and Yasser Al-Dosari. The external set contains 399 recordings (57 per reciter): four known reciters and three unseen reciters. Development data cover Surahs 78--103 and external data cover Surahs 104--114.

Public machine-readable per-recording metadata and source mappings are provided in the project repository.\footnote{\url{https://github.com/xyzeez/quran-reciter-classifier}} The CSV files record filename, reciter, augmentation status, duration, sample rate, run identifier, and feature index; filenames encode surah and ayah. Audio is not redistributed because source permissions may vary. Since each reciter is associated with one EveryAyah archive, surah separation prevents verse overlap but cannot exclude archive-, codec-, channel-, or session-specific cues.

The final development preprocessing record reports 2,535 successfully processed files, no skipped files, no augmented files, and 11,223.37 seconds of audio. The external preprocessing record reports 399 successfully processed files, no skipped files, no augmented files, and 1,724.84 seconds of audio. Each external reciter contributes exactly 57 recordings, preventing class-count imbalance in the external comparison. Clip durations remain variable after silence removal. Across the external set, they range from 1.46 to 11.24 seconds (median 4.19, mean 4.32); reciter-level medians range from 3.52 seconds for Maher Al Meaqli to 4.86 seconds for Mohammed Siddiq Al-Minshawi. Consequently, the reported results characterize verse-level recordings from the configured archives rather than fixed-duration studio speech.

% -----------------------------
% Table 1
% -----------------------------
\begin{table*}[t]
\centering
\caption{Dataset partitions used in the experimental evaluation.}
\label{table1}
\renewcommand{\arraystretch}{1.2}
\setlength{\tabcolsep}{6pt}
\resizebox{\textwidth}{!}{
\begin{tabular}{p{3cm} p{2.5cm} p{2cm} p{6.5cm}}
\toprule
\textbf{Partition} & \textbf{Recordings} & \textbf{Reciters} & \textbf{Details} \\
\midrule
Model-fitting set & 2,028 & 5 & Stratified 80\% of development corpus; Surahs 78--103 \\
Holdout set & 507 & 5 & Stratified 20\% of the same content/source distribution \\
External test set & 399 & 7 & Surahs 104--114; includes seen and unseen reciters \\
\bottomrule
\end{tabular}
}
\end{table*}

\subsection{Model Training Configuration}

Both models are trained on identical partitions, although their input representations differ. Table~\ref{table2} summarizes the training settings.

% -----------------------------
% Table 2
% -----------------------------
\begin{table}[htbp]
\centering
\caption{Model training configurations for Random Forest and BLSTM classifiers.}
\label{table2}
\renewcommand{\arraystretch}{1.2}
\setlength{\tabcolsep}{4pt}
\resizebox{\columnwidth}{!}{
\begin{tabular}{p{1.4cm} p{2.8cm} p{2.8cm}}
\toprule
\textbf{Parameter} & \textbf{Random Forest} & \textbf{BLSTM} \\
\midrule
Input features & 511-dim hybrid vector & 13 MFCC summaries repeated $\times 16$ \\
Architecture & 100 trees, depth 10 & BiLSTM (64 units) + attention \\
Optimizer & --- & AdamW (lr = 0.0005) \\
Batch size & --- & 64 \\
Regularization & Probability calibration & Dropout 0.5, early stopping \\
Random seed & 42 & 42 \\
\bottomrule
\end{tabular}
}
\end{table}

\subsection{Evaluation Protocol}

Evaluation uses four stages:

\begin{enumerate}
\item \textbf{Holdout Evaluation:} A stratified 20\% subset of the training dataset is used to evaluate model performance on unseen samples from known reciters.
\item \textbf{Closed-Set External Evaluation:} Evaluation is performed on 228 recordings from known reciters reciting previously unseen surahs (Surahs 104--114).
\item \textbf{Open-Set Evaluation:} Evaluation is performed on 171 recordings from unseen reciters not included in the training dataset.
\item \textbf{Combined External Evaluation:} A full evaluation on all 399 external test samples, combining both closed-set and open-set scenarios.
\end{enumerate}

This protocol separates model selection from final external assessment. The holdout set provides an internal estimate of known-reciter performance, while the external set is reserved for content-separated evaluation. For open-set testing, unknown-reciter recordings are not used as training classes; they are used only to test whether the reliability mechanism can reject out-of-distribution reciters.

\subsection{Evaluation Metrics}

Metrics include accuracy, macro F1, unseen-reciter rejection accuracy, false-positive identification rate (unknown accepted as known), and known-reciter false-rejection rate (known rejected as unknown). Wilson 95\% confidence intervals (CI) are reported for proportions. For open-set testing, an unseen-reciter sample is correct only when rejected as \textbf{Unknown}.

\subsection{Experimental Reproducibility}

Dataset splitting and RF training use random seed 42. Preprocessing, trained models, and per-file evaluation outputs are stored under versioned run identifiers. The reported BLSTM result is one run; deterministic CUDA settings and multi-seed variance were not evaluated.

Table~\ref{table_runs} identifies the exact saved artifacts used for the reported results. Both classifiers use development features from the same preprocessing run and are tested against the same external feature matrix. The detailed result files retain filename, duration, true label, known/unseen status, raw predicted label, final post-rejection label, confidence, reliability decision, normalized distance, and rejection reasons. These records permit the post-hoc count and confidence-interval analyses in this paper without recomputing features or model outputs.

\begin{table}[htbp]
\centering
\caption{Versioned artifacts underlying the reported evaluation.}
\label{table_runs}
\renewcommand{\arraystretch}{1.08}
\setlength{\tabcolsep}{3pt}
\resizebox{\columnwidth}{!}{
\begin{tabular}{l l}
\toprule
\textbf{Artifact} & \textbf{Run identifier} \\
\midrule
Development preprocessing & 20250526\_204349\_preprocess \\
External preprocessing & 20250525\_231932\_preprocess \\
RF model & 20250526\_205455\_train \\
BLSTM model & 20250526\_205529\_train \\
RF external evaluation & 20260616\_154714\_test \\
BLSTM external evaluation & 20260616\_154909\_test \\
\bottomrule
\end{tabular}}
\end{table}

\section{Results and Discussion}

On the 507-recording holdout, RF achieved 99.61\% (95\% CI: 98.57--99.89) and BLSTM 99.01\% (97.71--99.58). Before rejection, RF classified all 228 known external recordings correctly (100.0\%; 98.34--100.0), while BLSTM classified 223/228 correctly (97.81\%; 94.97--99.06). After applying Eq.~\ref{eq:reliability}, RF rejected 15/228 known recordings (6.58\%; 4.03--10.57) and BLSTM rejected 21/228 (9.21\%; 6.10--13.67). For 171 unseen-reciter recordings, RF rejected 167 (97.66\%; 94.14--99.09), while BLSTM rejected 149 (87.13\%; 81.29--91.35). Overall post-rejection accuracy was 95.24\% for RF and 88.97\% for BLSTM.

% -----------------------------
% Table 3
% -----------------------------
\begin{table}[htbp]
\centering
\caption{Summary of evaluation results across all experimental settings.}
\label{table3}
\renewcommand{\arraystretch}{1.1}
\setlength{\tabcolsep}{3pt}
\resizebox{\columnwidth}{!}{
\begin{tabular}{l c c}
\toprule
\textbf{Evaluation setting} & \textbf{RF} & \textbf{BLSTM} \\
\midrule
Holdout accuracy ($n=507$) & 99.61\% & 99.01\% \\
Known raw accuracy ($n=228$) & 100.0\% & 97.81\% \\
Known FRR ($n=228$) & 6.58\% & 9.21\% \\
Unseen rejection ($n=171$) & 97.66\% & 87.13\% \\
Unseen FPR ($n=171$) & 2.34\% & 12.87\% \\
Overall accuracy ($n=399$) & 95.24\% & 88.97\% \\
\bottomrule
\end{tabular}
}
\end{table}

\subsection{Class-Level Holdout Behavior}

The stored classification reports permit a more granular audit of the random holdout. Table~\ref{table_holdout_f1} reports F1 scores rather than accuracy alone. Both models attain perfect F1 for Abdulrahman Alsudaes and Ahmad Al-Ajmy. RF's lowest class F1 is 0.9902 for Yasser Al-Dosari, while BLSTM's lowest is 0.9758 for Maher Al Meaqli; BLSTM also has lower recall for Abdulbasit Abdulsamad (0.9706 versus 0.9902 for RF). The corresponding macro-F1 scores are 0.9961 and 0.9902. These high values characterize the in-distribution holdout only: because it is sampled from the same processed corpus and surah range as training, it is less demanding than the content-separated external test.

\begin{table}[htbp]
\centering
\caption{Per-class F1 on the 507-recording random holdout.}
\label{table_holdout_f1}
\renewcommand{\arraystretch}{1.06}
\setlength{\tabcolsep}{4pt}
\begin{tabular}{l r r r}
\toprule
\textbf{Reciter} & \textbf{$n$} & \textbf{RF} & \textbf{BLSTM} \\
\midrule
Abdulbasit Abdulsamad & 102 & 0.9951 & 0.9851 \\
Abdulrahman Alsudaes & 101 & 1.0000 & 1.0000 \\
Ahmad Al-Ajmy & 101 & 1.0000 & 1.0000 \\
Maher Al Meaqli & 102 & 0.9951 & 0.9758 \\
Yasser Al-Dosari & 101 & 0.9902 & 0.9901 \\
\midrule
Macro average & 507 & 0.9961 & 0.9902 \\
\bottomrule
\end{tabular}
\end{table}

\subsection{Per-Reciter Rejection Analysis}

Rejection behavior varies by reciter (57 recordings each). For the unseen Abdullah Al-Johany, Mohammed Siddiq Al-Minshawi, and Saad Al-Ghamdi, RF rejection is 100.0\%, 92.98\%, and 100.0\%, while BLSTM rejection is 100.0\%, 77.19\%, and 84.21\%. For the known Abdulbasit Abdulsamad, Abdulrahman Alsudaes, Maher Al Meaqli, and Yasser Al-Dosari, RF false rejection is 3.51\%, 8.77\%, 10.53\%, and 3.51\%; BLSTM false rejection is 3.51\%, 7.02\%, 12.28\%, and 14.04\%.

Table~\ref{table_errors} separates raw classifier correctness from rejection behavior. For known reciters, ``Raw correct'' counts label decisions before the reliability rule, ``Reject'' counts known inputs changed to \textbf{Unknown}, and ``Accepted error'' counts incorrect known labels that passed the rule. For unseen reciters, raw correctness is undefined because no unseen identity is a training class; accepted errors are false-positive identifications. RF classifies every known external recording correctly before rejection, but rejects 15 of them. BLSTM makes five raw known-reciter errors for Yasser Al-Dosari; four are rejected and one is accepted as an incorrect known identity. The 26 unseen false acceptances across both models are concentrated in two reciters: Al-Minshawi for both models and Saad Al-Ghamdi for BLSTM.

\begin{table}[htbp]
\centering
\caption{Per-reciter decisions at the fixed operating point (57 recordings each). Entries are raw correct/rejected/accepted error.}
\label{table_errors}
\renewcommand{\arraystretch}{1.08}
\setlength{\tabcolsep}{3pt}
\resizebox{\columnwidth}{!}{
\begin{tabular}{l c c c}
\toprule
\textbf{Reciter} & \textbf{Status} & \textbf{RF C/R/E} & \textbf{BLSTM C/R/E} \\
\midrule
Abdulbasit Abdulsamad & Known & 57/2/0 & 57/2/0 \\
Abdulrahman Alsudaes & Known & 57/5/0 & 57/4/0 \\
Maher Al Meaqli & Known & 57/6/0 & 57/7/0 \\
Yasser Al-Dosari & Known & 57/2/0 & 52/8/1 \\
Abdullah Al-Johany & Unseen & ---/57/0 & ---/57/0 \\
Mohammed Siddiq Al-Minshawi & Unseen & ---/53/4 & ---/44/13 \\
Saad Al-Ghamdi & Unseen & ---/57/0 & ---/48/9 \\
\bottomrule
\end{tabular}}
\end{table}

\subsection{Coverage and Selective Accuracy}

Open-set rejection converts the classifier into a selective predictor: it returns an identity only for inputs satisfying all reliability conditions. Table~\ref{table_selective} therefore reports both coverage (the fraction receiving a known identity) and accuracy conditional on acceptance. RF accepts 217/399 inputs and correctly labels 213 of them, giving 54.39\% coverage and 98.16\% selective accuracy. BLSTM accepts 229/399 and correctly labels 206, giving 57.39\% coverage but 89.96\% selective accuracy. Among accepted known-reciter recordings, RF is correct on 213/213 and BLSTM on 206/207. The larger BLSTM coverage is thus not an unqualified benefit: it includes 22 accepted unseen-reciter inputs, compared with four for RF. This view complements overall accuracy by separating abstention frequency from the trustworthiness of non-abstained predictions.

\begin{table}[htbp]
\centering
\caption{Selective-prediction behavior on the 399 external recordings.}
\label{table_selective}
\renewcommand{\arraystretch}{1.08}
\setlength{\tabcolsep}{3pt}
\resizebox{\columnwidth}{!}{
\begin{tabular}{l r r r r}
\toprule
\textbf{Model} & \textbf{Accepted} & \textbf{Coverage} & \textbf{Selective acc.} & \textbf{Unseen accepted} \\
\midrule
RF & 217/399 & 54.39\% & 213/217 (98.16\%) & 4/171 \\
BLSTM & 229/399 & 57.39\% & 206/229 (89.96\%) & 22/171 \\
\bottomrule
\end{tabular}}
\end{table}

\subsection{False-Acceptance Audit}

The per-file records show that RF's four unseen false acceptances all come from Al-Minshawi: three are labeled Abdulbasit Abdulsamad and one Abdulrahman Alsudaes. Their accepted confidences range from 0.9543 to 0.9800 and their normalized distance ratios from 0.2940 to 0.9511, so each satisfies the fixed rule despite being an unseen identity. BLSTM's 22 false acceptances comprise 13 Al-Minshawi and nine Saad Al-Ghamdi recordings. Twelve are labeled Abdulrahman Alsudaes, eight Maher Al Meaqli, and one each Abdulbasit Abdulsamad and Ahmad Al-Ajmy; confidence ranges from 0.9526 to 0.9993 and distance ratio from 0.1482 to 0.9599. Thus, the remaining errors are not merely low-confidence boundary cases. They are confidently accepted samples inside the model-specific distance limits, motivating calibrated threshold sweeps and channel-controlled evaluation rather than a claim that the present fixed rule fully solves unknown-reciter detection.

\subsection{Reliability-Criterion Diagnostics}

The saved rejection reasons also permit a criterion-level audit at the selected thresholds. Table~\ref{table_failures} counts how often each condition fails across all 399 external recordings. Counts overlap because one recording can fail several conditions. RF produces more confidence-related failures than BLSTM, whereas the distance condition fails for almost the same number of recordings (119 and 120). This helps explain why BLSTM accepts more unseen inputs: its posterior outputs fail the primary-confidence, secondary-confidence, and margin requirements less often, even though the shared distance representation produces a similar number of distance failures. These counts describe one operating point and should not be interpreted as a threshold sweep.

\begin{table}[htbp]
\centering
\caption{Failure occurrences for reliability conditions ($n=399$; non-exclusive counts).}
\label{table_failures}
\renewcommand{\arraystretch}{1.08}
\setlength{\tabcolsep}{4pt}
\begin{tabular}{l r r}
\toprule
\textbf{Failed condition} & \textbf{RF} & \textbf{BLSTM} \\
\midrule
$p_1 < 0.95$ & 162 & 115 \\
$p_2 \geq 0.10$ & 117 & 78 \\
$p_1-p_2 < 0.80$ & 140 & 84 \\
$d(x,c_1)/T(c_1) > 1.0$ & 119 & 120 \\
\bottomrule
\end{tabular}
\end{table}

\subsection{Discussion}

At the selected operating point, the hybrid-feature RF pipeline outperforms the MFCC-summary BLSTM baseline, but unequal representations preclude a general RF-versus-deep-learning conclusion. Rejection also trades lower unseen acceptance for false rejection of known reciters.

\subsection{Threats to Validity}

This proof of concept uses five development reciters and only three unseen reciters from one provider. It does not control microphone, environment, riwayah, archive, codec, channel, or recording session; source characteristics may therefore correlate with identity. The random holdout shares content range and sources with model fitting. Surah separation evaluates new content but cannot prove purely reciter-specific modeling. The RF and BLSTM inputs are unequal, thresholds are manually set, only one reported run is available, and no threshold-independent AUROC/OSCR or sensitivity analysis was performed. Confidence intervals quantify sampling uncertainty for the fixed test sets, not training-seed variance.

\subsection{Future Work}

Future work should expand reciters, riwayat, microphones, environments, and channel-separated sessions; evaluate genuine frame-level, feature-matched sequence models across seeds; and report threshold sweeps, AUROC, and OSCR.

\section{Conclusion}

This paper presented a small-scale open-set Qur'anic reciter recognition proof of concept with content-separated external evaluation. At the selected rejection operating point, hybrid-feature RF achieved 97.66\% unseen-reciter rejection and 95.24\% overall external accuracy, compared with 87.13\% and 88.97\% for the MFCC-summary BLSTM baseline. The corresponding known-reciter false-rejection rates were 6.58\% and 9.21\%. These findings support the implemented hybrid-feature RF pipeline on this corpus, while the unequal inputs, limited reciter set, single-provider recordings, and manual thresholds prevent a general architecture-level conclusion. Broader, channel-controlled and feature-matched evaluation is required before deployment claims can be made.

\bibliographystyle{IEEEtran}
\bibliography{references}

\end{document}
